\documentclass[preprint,12pt]{elsarticle}

\usepackage{fontspec}
\usepackage{xeCJK}
\usepackage{amsmath,amssymb}
\usepackage{booktabs}
\usepackage{longtable}
\usepackage{array}
\usepackage{multirow}
\usepackage{fvextra}
\usepackage{graphicx}
\usepackage{float}
\usepackage{hyperref}
\usepackage{pdflscape}
\usepackage{xcolor}

\setmainfont{Times New Roman}
\setsansfont{Helvetica}
\setmonofont{Menlo}
\setCJKmainfont[AutoFakeBold=3,AutoFakeSlant=.2]{Noto Serif CJK SC}
\setCJKsansfont[AutoFakeBold=3]{Noto Sans CJK SC}
\setCJKmonofont{Sarasa Mono SC}

\hypersetup{
  unicode=true,
  colorlinks=true,
  linkcolor=blue!50!black,
  citecolor=blue!50!black,
  urlcolor=blue!50!black
}

\journal{}

\setlength{\LTleft}{0pt}
\setlength{\LTright}{0pt}
\setlength{\LTpre}{6pt}
\setlength{\LTpost}{6pt}
\setlength{\tabcolsep}{4pt}
\setlength{\emergencystretch}{4em}
\renewcommand{\arraystretch}{1.05}
\sloppy

\providecommand{\tightlist}{}
\DefineVerbatimEnvironment{verbatim}{Verbatim}{breaklines=true,breakanywhere=true,fontsize=\small}
\AtBeginEnvironment{longtable}{\footnotesize}
\newsavebox{\pandocbox}
\newcommand{\pandocbounded}[1]{%
  \sbox{\pandocbox}{#1}%
  \ifdim\wd\pandocbox>\linewidth
    \resizebox{\linewidth}{!}{\usebox{\pandocbox}}%
  \else
    \usebox{\pandocbox}%
  \fi
}

\begin{document}

\begin{frontmatter}

\title{数据要素利用与资本市场定价效率\\
\large 基于年报文本分析与大语言模型语义评分的证据}

\author{匿名作者}

\begin{abstract}
年报中的数据要素利用信息是观察企业数字实践如何进入资本市场价值发现的重要窗口。本文基于2011至2024年沪深A股上市公司年报全文，利用关键词文本分析与大语言模型语义评分构建双测度体系，考察年报中的数据要素利用信息对资本市场定价效率的影响。其中，DU\_kw刻画年报中的数据要素利用信息广度，是本文的基准处理变量；DU\_llm及其长度标准化版本DU\_llm\_lenstd刻画语义深度，用于检验深度信息是否支持同方向结论。进一步地，本文基于五维关键词词典和年报章节信息构造了数据利用链条完整度、闭环式表述和MD\&A章节限定等增强测度，以检验测度能否从泛化数字叙事进一步推进到更接近数据资源实际使用的口径。研究发现，年报中的数据要素利用信息显著降低股价延迟。该结论在滞后规格、工具变量、双重机器学习、Oster界检验、Heckman两阶段法和安慰剂检验下保持稳定，多种识别练习支持因果解释。构念边界检验显示，在控制年报篇幅和剔除八个最泛化的数字叙事词后，核心关系仍然成立；但与GenericNarr的联合回归表明，相关披露与更宽泛的数字叙事仍存在交叠，因此对概念独立性的表述应保持审慎。增强测度检验进一步表明，更完整的数据利用链条、闭环式表述以及MD\&A中的相关披露同样与更低的股价延迟相关，但这些增强指标在与DU\_kw或GenericNarr联合进入时通常不能稳定替代基准处理变量，因此更适合作为支持性证据而非替换口径。机制检验表明，年报中的数据要素利用信息与更低的现金流波动和更低的供应链集中度稳健相关，并在分析师预测分歧维度上呈现补充性证据。扩展分析进一步显示，关键词广度与语义深度的错配会加剧股价延迟，而语义深度在与关键词广度同时进入时仍保留额外解释力。研究结论有助于理解年报中的数据要素利用信息的资本市场价值，并为完善相关披露与估值指引提供经验证据。
\end{abstract}

\begin{highlights}
\item 基于上市公司年报全文，构建关键词频率与大语言模型语义评分的双测度体系，并进一步报告闭环度、MD\&A限定和长度标准化语义深度等增强测度。
\item 年报中的数据要素利用信息显著降低股价延迟；滞后规格、工具变量、双重机器学习与多项辅助识别练习均支持这一结论。
\item 更完整的数据利用链条、闭环式表述以及MD\&A中的相关披露同样与更低股价延迟相关，但这些增强口径更适合作为支持性证据而非替换基准处理变量。
\item 机制结果显示，年报中的数据要素利用信息与更低的现金流波动和更低的供应链集中度稳健相关，并在分析师预测分歧维度上呈现补充性证据。
\item 构念边界检验表明，核心关系在篇幅控制和去泛化词口径下仍然成立，但与更宽泛数字叙事的边界并非完全分离，概念独立性应作审慎表述。
\end{highlights}

\begin{keyword}
数据要素利用 \sep 资本市场定价效率 \sep 股价延迟 \sep 年报文本分析 \sep 大语言模型
\end{keyword}

\end{frontmatter}

\section{一、引言}\label{ux4e00ux5f15ux8a00}

数据要素在企业价值创造中的作用持续上升，但会计计量体系仍然滞后（洪永淼和史九领，2024）。2022年``数据二十条''和2024年《企业数据资源相关会计处理暂行规定》在制度层面迈出了重要一步，然而现实困难依然突出：历史成本模式难以反映数据资源的动态价值特征，商业模式差异又使同一数据资源在不同场景下价值迥异（黄世忠等，2023；Mohan等，2026）。这一计量缺口意味着投资者难以仅凭财务报表完整评估企业价值，资本市场对相关企业的定价效率因此受到影响。

企业在年报中披露的数据要素利用信息可能通过两条相反的路径影响定价效率。从信息供给的角度看，数据利用的实质性描述降低了投资者面临的信息不对称，有助于加速价格发现（Diamond和Verrecchia，1991）。但从信息模糊的角度看，数据驱动商业模式的内在复杂性和估值标准的缺失可能加剧不确定性，反而阻碍价格发现（Zhang，2006）。净效果取决于哪种力量占主导，需要经验数据加以识别。中国A股市场的年报经过审计、受信息披露监管约束，数据要素政策框架也在逐步完善，这些条件可能有利于信息供给效应的发挥，但仍需实证检验。

既有研究主要沿两条脉络展开。第一条脉络关注数据要素对企业行为的影响。朱康和唐勇（2025）发现数据要素利用降低了金融资产配置并促进实体投资，李姝等（2025）提供了数据资产缓解融资约束的互补证据，Shao等（2024）从经济政策不确定性视角报告了类似发现。这些研究揭示了数据利用在企业决策层面的实质效应，但研究视角大多停留在企业行为层面，尚未系统检验资本市场如何吸收这类信息。

第二条脉络直接考察信息环境与定价效率的关系。会计信息质量的改善能够降低股价延迟（Callen等，2013），非财务信息的前瞻性有助于提升定价效率（Zhang和Wang，2024），另类数据对财务预测的改善在较长时间跨度上更为明显（Dessaint等，2024）。在中国市场，行业信息披露指引（Gao等，2024）和主流财经媒体关注（Yin和Hu，2024）均被发现能改善股价信息含量。与数据要素更直接相关的研究中，数字化转型降低股价延迟的初步证据已被报告（Feng等，2024），数据资产信息披露降低股价同步性和股价延迟的结论也在特定样本中得到支持（Sun和Du，2024；Chen，2024；Huang等，2026；Wei等，2025；李世刚等，2025）。

然而，``数字化转型''\,``数据资产披露''与``数据要素利用''是相关但不等同的概念。数字化转型侧重技术采纳，数据资产披露侧重会计确认，数据要素利用关注的是企业对数据资源的实际应用。本文后续的构念边界检验也表明，这三者并非可以被完全机械切分：数据利用测度在剔除泛化词和控制篇幅后仍保留解释力，但与更宽泛数字叙事之间仍存在交叠。因此，本文的目标不是将其表述为与数字化转型完全正交的概念，而是更审慎地识别其相对更贴近``数据资源实际利用''的部分。

如果年报中的数据利用信息确实帮助市场更快理解企业价值，那么这种作用不应只体现在总量上的股价延迟下降，还应体现在更细的传导环节上。其一，信息披露可能降低分析师对企业基本面的认知分歧，使预测更趋一致。其二，数据驱动的经营优化可能降低企业现金流波动，从而减少估值不确定性。其三，数据能力可能帮助企业优化客户和供应商匹配，缓解供应链集中度过高带来的基本面脆弱性。这三条路径分别对应信息质量、经营稳定性和供应链韧性，能够帮助识别数据要素利用如何从年报文本传导至股票价格。

本文旨在构建这一完整路径。与朱康和唐勇（2025）聚焦企业决策行为和李世刚等（2025）以入表口径度量数据资产化不同，本文直接考察年报中数据要素实际应用的实质性描述如何影响股价的信息吸收速度，并系统检验其与分析师预测分歧、现金流波动和供应链集中度之间的关系。

本文以2011至2024年沪深A股上市公司为样本，基于年报全文考察数据要素利用信息对股价延迟的影响。度量上采用两种互补策略：关键词频率指标扫描年报中与数据收集、分析、应用和价值变现相关的表述，捕捉利用信息的广度；大语言模型对关键词上下文进行语义评分，刻画利用信息的深度。本文将DU\_kw作为基准处理变量，将DU\_llm作为语义增强测度，并进一步构造长度标准化的DU\_llm\_lenstd用于检验语义深度是否在控制篇幅后仍保留解释力。与此同时，本文进一步利用现有五维词典和年报章节特征，构造数据利用链条完整度（DUchain\_count）、闭环式表述（DUclosedloop）、窄口径核心强度（DUcore）和MD\&A章节限定口径（DUkw\_mda），以检验测度能否从泛化数字叙事进一步推进到更接近数据资源实际使用的口径。研究发现，年报中的数据要素利用信息显著降低股价延迟；进一步以前一期DU解释当期PriceDelay，并结合工具变量、双重机器学习、Oster界检验、Heckman两阶段法和安慰剂检验，结论保持稳定，多种识别练习支持因果解释。构念边界检验显示，核心结果在篇幅控制和去泛化词口径下仍然成立，但GenericNarr联合回归表明二者存在较强重叠，因此概念独立性应作审慎表述。增强测度结果显示，更完整的数据利用链条、闭环式表述以及MD\&A中的相关披露同样与更低的股价延迟相关，但这些口径在与DU\_kw或GenericNarr联合进入时通常不能稳定替代基准处理变量，因此更适合作为支持性测度。机制结果显示，数据要素利用信息与更低的现金流波动和更低的供应链集中度稳健相关，并在分析师预测分歧维度上呈现补充性证据。扩展分析进一步表明，关键词广度与语义深度的错配会加剧股价延迟，而语义深度在与关键词广度同时进入时仍保留额外解释力。异质性结果进一步表明，该效应不仅在2022年之前更强，也在高分析师关注企业和非主板企业中表现得更为明显；此外，行业整体披露水平对企业股价延迟存在独立的负向作用，绝对量级在滞后规格下甚至超过企业自身披露，体现出行业层面的信息环境同群溢出。

本文的贡献主要体现在以下三个方面。第一，构建了关键词频率与大语言模型语义评分的双测度体系，并进一步通过去泛化词、联合回归、五维闭环度、闭环式表述和MD\&A章节限定口径检验测度边界。本文据此不再将``数据要素利用''表述为与更宽泛数字叙事完全分离的概念，而是更审慎地识别其相对更贴近实际数据利用的部分；其中，闭环度和经营章节口径为``更完整、更嵌入''的数据利用信息提供了支持性证据。第二，本文将数据要素研究从企业行为后果推进到资本市场的信息效率后果，并从信息分歧、经营稳定性与供应链韧性三个维度识别传导路径；进一步的WashGap分析显示，市场更倾向于奖励具有实质内容和经营嵌入的数据利用披露，而非单纯广泛但浅层的概念性表述。第三，本文将时间顺序更清晰的滞后规格前置，并结合工具变量、双重机器学习、Oster界检验、Heckman选择修正和安慰剂检验，形成多层次的识别框架，以更审慎的方式支持因果解释。

本文余下部分安排如下。第二部分建立理论框架并导出研究假说。第三部分介绍研究设计。第四部分依次报告基准回归、构念边界与测度补强、时间顺序与内生性识别、双重机器学习估计和稳健性检验。第五部分进一步检验传导机制与广度---深度错配。第六部分报告政策阶段、信息中介活跃度和上市板块异质性分析。第七部分总结全文并讨论政策启示。

\section{二、理论分析与研究假说}\label{ux4e8cux7406ux8bbaux5206ux6790ux4e0eux7814ux7a76ux5047ux8bf4}

\subsection{（一）数据要素利用与资本市场定价效率}\label{ux4e00ux6570ux636eux8981ux7d20ux5229ux7528ux4e0eux8d44ux672cux5e02ux573aux5b9aux4ef7ux6548ux7387}

在剩余收益估值框架下，企业经济价值以账面净资产为锚（Ohlson，1995）。数据资产的大规模表外存在动摇了这一锚点：大量数据资源因不满足可辨认性或可靠计量条件而停留在报表之外，现行成本法又无法反映数据资源随应用场景扩展而非线性增值的特征（黄世忠等，2023），账面价值因此系统性低估了企业的真实经济价值。无形资产的定价偏差在近年日趋突出，中国A股市场的偏差尤为明显（Chen, J.等，2024；Dong和Doukas，2025；Bongaerts等，2025）。2024年施行的入表新规虽将数据资源纳入资产负债表，但适应期内投资者解读这类信息的难度并未立即下降。

这一计量缺口使市场参与者难以仅凭财务报表评估企业真实价值。企业在年度报告中披露的数据要素利用信息因此成为弥补计量缺口的增量信号，但这类信号对定价效率的影响方向并不确定。

信息供给效应强调信息不对称的改善作用。可信的公开信息披露压缩了知情与非知情交易者之间的信息差异，加速价格发现过程，这一机制既有理论上的严格论证（Grossman和Stiglitz，1980；Diamond和Verrecchia，1991），也有中国市场行业信息披露指引改善股价信息含量的独立实证支持（Gao等，2024）。信息模糊效应则指向相反方向，即信息不确定性越高，投资者越难以准确定价，价格发现因此越迟缓（Zhang，2006），数据驱动商业模式的内在复杂性、会计准则的估值空白以及数据丰裕条件下信息处理能力的制约（Dugast和Foucault，2025）都可能加剧这种不确定性。两种效应的行为假设不同，前者认为价格发现的约束在于信息供给不足，后者认为瓶颈在于投资者理解和处理信息的能力，净结果取决于哪种力量占主导。

在中国A股市场，若干制度条件可能有利于信息供给效应的发挥。年报是经过审计的法定文件，数据要素利用描述受到会计准则和审计程序的约束，其可验证性高于自愿性渠道。2022年``数据二十条''和2024年入表新规逐步降低了理解和评估数据资产信息的门槛，制度框架的渐进完善使信息模糊效应的制约有所减弱。若信息供给效应占优，应观察到企业数据要素利用与股价延迟呈显著负向关联。据此提出：

H1：企业数据要素利用降低了资本市场的股价延迟，即提升了定价效率。

\subsection{（二）定价效率的传导机制}\label{ux4e8cux5b9aux4ef7ux6548ux7387ux7684ux4f20ux5bfcux673aux5236}

上述假说关注的是总效应，但信息从企业年报向股票价格的具体传导路径尚待厘清。本文重点检验三条机制路径。

第一条是信息共识渠道。数据要素利用涉及技术架构、业务流程和价值实现方式，具有较强的专业门槛。如果年报披露能够提升外部投资者对企业数据利用方式的理解，分析师在处理相关信息时应当形成更一致的判断，预测分歧随之下降。分析师分歧的收敛意味着信息解释框架更趋一致，进而有助于相关信息更快融入股价。据此提出：

H2a：企业数据要素利用降低了分析师预测分歧。

第二条是经营稳定性渠道。数据驱动的运营管理有助于提升需求预测、库存配置和资源调度的效率，使企业经营活动更具可预测性。若数据要素利用缓解了经营层面的随机波动，企业现金流应更平稳，投资者面对的基本面不确定性也会下降，价格发现过程因而更顺畅。据此提出：

H2b：企业数据要素利用降低了现金流波动率。

第三条是供应链韧性渠道。供应链集中度过高意味着企业对少数关键客户或供应商依赖较强，一旦交易对手发生变化，基本面冲击更容易放大。数据要素利用可以帮助企业更准确地识别、筛选和匹配潜在交易伙伴，提高供应链管理效率并促进交易关系多元化。若这一过程成立，企业的综合供应链集中度应下降，基本面风险的暴露也将减弱。据此提出：

H2c：企业数据要素利用降低了综合供应链集中度。

\subsection{（三）效应的异质性}\label{ux4e09ux6548ux5e94ux7684ux5f02ux8d28ux6027}

总效应和传导机制之外，制度环境差异也可能调节定价效应的强弱。

第一，政策阶段差异影响数据利用信息的边际新颖性。相较于2020或2021，2022年``数据二十条''更接近数据要素制度密集推进的起点，也更适合作为主文中的阶段分界。在较早阶段，数据利用信息对投资者而言更具新颖性和辨识度，其边际定价价值可能更高；随着披露行为趋于制度化和常态化，边际信息增量可能减弱。据此提出：

H3a：企业数据要素利用对股价延迟的改善效应在2022年之前更强。

第二，信息中介活跃度可能影响数据利用披露的吸收速度。数据要素利用表述具有一定专业门槛，若分析师等信息中介对企业持续跟踪更充分，年报中的相关信息更容易被解释、传播并进入价格。相反，在分析师关注较低的企业中，即便存在披露，信息也可能因处理能力不足而更慢反映到价格中。据此提出：

H3b：企业数据要素利用对股价延迟的改善效应在高分析师关注企业中更强。

第三，上市板块差异可能改变数据利用披露的边际价值。主板公司通常经营模式更成熟、投资者覆盖更广、常规信息供给更充分，相关披露未必形成强增量信号；而非主板公司往往具有更强的成长性和更高的信息摩擦，数据利用披露可能更容易成为投资者重新理解企业价值的重要线索。据此提出：

H3c：企业数据要素利用对股价延迟的改善效应在非主板企业中更强。

产业属性差异同样可能影响市场的先验预期。例如，数字经济核心产业的企业天然与数据技术联系更紧密，相关披露未必构成强增量信息。但在统一滞后规格下，该类产业差异并未得到稳定支持，因此本文仅将其作为附录中的补充异质性结果报告，而不再据此建立主文的正式假说。

\section{三、研究设计}\label{ux4e09ux7814ux7a76ux8bbeux8ba1}

\subsection{（一）样本选择与数据来源}\label{ux4e00ux6837ux672cux9009ux62e9ux4e0eux6570ux636eux6765ux6e90}

本文以2011至2024年沪深A股上市公司为样本。数据包括两个层面，年报全文用于提取数据要素利用指标，CSMAR数据库的财务、治理和市场数据用于构造因变量和控制变量。样本筛选依据中国证监会行业分类标准剔除J类金融保险业企业，剔除ST/*ST企业，剔除上市时间在观测年份内不足12个月的企业-年份观测。对连续变量采取1\%和99\%分位数的Winsorize缩尾处理。

变量的时间匹配遵循按报告期对齐的原则（Feng等，2024；Wei等，2025）：t年年报覆盖企业t年经营活动，PriceDelay基于t年日收益率计算，二者同属一个会计年度。考虑到英文审稿中可能提出的``用事后文本编码当年信息环境''的质疑，本文在第四部分除报告同期规格外，还将\(t-1\)期DU解释\(t\)期PriceDelay的滞后规格前置呈现，作为时间顺序更清晰的偏好识别规格。年报文本提取的时间范围为2010至2024年，其中2010年数据仅用于该滞后分析。

经过上述筛选，最终获得43,735个企业-年份观测，涵盖约4,950家独立企业。OLS基准回归和工具变量估计均基于该样本，扣除固定效应单例后进入回归的观测约43,700。DML估计因使用27个控制变量和交叉拟合，样本约35,700，在相应章节单独说明。

年报全文来源于上市公司向证监会提交的法定年报TXT格式文件。因变量PriceDelay基于CSMAR数据库的股票日收益率数据构造，控制变量和机制变量同样来自CSMAR。

\subsection{（二）数据要素利用的测度}\label{ux4e8cux6570ux636eux8981ux7d20ux5229ux7528ux7684ux6d4bux5ea6}

数据资产及其相关活动长期处于会计体系之外，传统财务指标无法直接观察（Coyle和Manley，2024），文本分析是目前获取相关信息的主要方式。本文沿数据价值链的不同环节构建了五维度关键词体系，共计73个关键词（见下表）。

\textbf{表1 数据要素利用关键词体系}

{\def\LTcaptype{} % do not increment counter
\begin{longtable}[]{@{}lll@{}}
\toprule\noalign{}
维度 & 关键词数 & 示例 \\
\midrule\noalign{}
\endhead
\bottomrule\noalign{}
\endlastfoot
数据存量 data\_stock & 15 & 大数据、数据库、数据中心 \\
数据开发能力 data\_dev & 19 & 数据挖掘、机器学习、数字化转型 \\
数据驱动应用 data\_app & 19 & 精准营销、智能推荐、风控模型 \\
数据价值变现 data\_value & 11 & 数据资产、数据交易、数据入表 \\
数据治理 data\_gov & 9 & 数据安全、数据隐私、数据合规 \\
\end{longtable}
}

这一分类方案参考既有数字化文本测度，并针对朱康（2025）采用的两维度体系进行了细化，新增数据开发能力、数据价值变现和数据治理三个维度，覆盖数据要素利用的主要环节。完整关键词清单见附录。

关键词频率的计算方法为：

\[DU_{kw} = \frac{\text{关键词总出现次数}}{\text{年报全文总字数}} \times 10000 \qquad (1)\]

即每万字中数据相关关键词的出现次数。文本处理采用多编码支持策略和长词优先匹配策略，以准确读取多种编码格式并避免短词误匹配。

关键词频率统计存在测量噪声：企业可能在年报中笼统地提及``大数据''\,``数字化''而并无实质利用，也可能在不同语境下使用相同关键词但利用深度迥异。为缓解这一问题，本文引入大语言模型语义评分作为第二套测度。

具体步骤如下：首先基于73个关键词定位年报中所有命中位置及其前后100字的上下文窗口；然后使用大语言模型对每个上下文进行0-3分的语义评分，判断该处关键词是否反映了实质性的数据要素利用（0=无实质内容，3=深度利用）；最后取企业-年份层面各上下文评分的众数作为该企业当年的LLM评分llm\_score，以抑制少数噪声上下文的拉动。评分使用Claude Haiku 4.5模型，温度参数设为0以确保输出确定性。附录进一步报告评分模板、0/1/2/3判分规则和真实上下文示例。

在此基础上构造三个衍生指标。第一，LLM二值化指标llm\_binary在llm\_score\textgreater=2时取1否则取0。第二，归档语义增强指标DU\_llm = ln(1 + kw\_total) × (llm\_score / 3)，将关键词总出现次数取对数后与归一化的语义质量相乘。第三，长度标准化语义深度指标DU\_llm\_lenstd = ln(1 + DU\_kw) × (llm\_score / 3)，以已做篇幅标准化的DU\_kw替换kw\_total，对语义深度进行更保守的口径校正。本文将DU\_kw明确作为基准处理变量和主要识别对象，将DU\_llm作为归档语义测度并行报告，并以DU\_llm\_lenstd检验语义深度在控制篇幅后的稳健性。

为识别更宽泛数字叙事与数据利用披露之间的边界，本文进一步构造GenericNarr，即八个最泛化数字叙事词------``数字化转型、数字化、智能化、信息化、人工智能、大数据、智能制造、智慧城市''------的每万字出现频次；并据此定义去泛化词指标DU\_kw\_strict，即从kw\_total中扣除上述八个词后的每万字频次。扩展分析还构造WashGap = z(DU\_kw) - z(DU\_llm\_lenstd)，刻画``提得多但语义深度不足''的广度---深度错配；同时定义BroadShallow虚拟变量，当企业当年DU\_kw处于年度上三分位且DU\_llm\_lenstd处于年度下三分位时取1，否则取0。

稳健性检验还使用两个备选指标：DU\_kw\_ln = ln(1 + DU\_kw)，即对频率进行对数变换以处理右偏分布；DU\_sub\_ln = ln(1 + DU\_sub)，其中DU\_sub为基于规则判断的实质利用次数。

在上述双测度之外，本文进一步构造四类增强口径，以检验测度能否从``泛化数字叙事''推进到``更接近数据资源实际使用''的表述。第一，DUevent\_raw和DUevent分别定义为规则识别的实质利用次数及其每万字标准化值。这里的规则识别基于关键词附近上下文中的动作词、具体业务对象和数值信息，substantive\_count越高，表示文本中越多表述具有``被用于具体经营或治理场景''的特征。第二，DUchain\_count定义为五个维度中被命中的维度数，取值为0至5；DUclosedloop为闭环式表述虚拟变量，当企业同时出现数据资源/存量类、应用/开发类以及治理/价值实现类表述时取1，否则取0。第三，DUcore\_raw和DUcore聚焦数据驱动应用、价值实现和治理合规三类更接近经营与制度闭环的词项强度，用于形成更窄的支持性口径。第四，DUkw\_mda为MD\&A章节内关键词频率，衡量数据利用表述是否更集中于经营讨论部分。本文将这些增强指标明确定位为支持性测度：其目的在于检验``更完整、更嵌入''的数据利用信息是否同样具有资本市场定价含量，而非替换DU\_kw这一基准处理变量。

指标有效性的验证如下。时间趋势上，DU\_kw和DU\_llm在2011至2024年间均呈上升趋势，在政策时点加速明显。行业分布上，信息技术产业和科学研究与技术服务业的指标均值远高于制造业和采矿业，符合产业间数据依赖程度的差异。两套测度的Pearson相关系数为0.678，一致性较强但并不重叠：DU\_kw侧重利用的广度，DU\_llm侧重利用的深度。本文进一步通过篇幅控制、去泛化词口径和GenericNarr联合回归检验构念边界，并增加DUchain\_count、DUclosedloop、DUcore和DUkw\_mda等增强口径，以检验更接近业务嵌入和链条完整性的表述是否同样能够预测股价延迟。需要说明的是，受限于现有归档条件，本文暂未能重建基于mean score和share(score\textgreater=2)的全样本替代聚合；但本文已对120条分层上下文样本完成独立人工复核，并报告人工评分与独立Codex样本级评分的一致性统计，以补强语义评分规则的可复核性。

\subsection{（三）变量定义}\label{ux4e09ux53d8ux91cfux5b9aux4e49}

因变量为股价延迟PriceDelay，采用Hou和Moskowitz（2005）的方法构造。对企业\(i\)在年份\(t\)的日收益率序列，要求至少120个交易日，分别估计受限模型和无限制模型：

受限模型： \[r_{i,d} = \alpha + \beta_0 \cdot r_{m,d} + \varepsilon_{i,d} \qquad (2)\]

无限制模型： \[r_{i,d} = \alpha + \beta_0 \cdot r_{m,d} + \sum_{k=1}^{5} \beta_k \cdot r_{m,d-k} + \varepsilon_{i,d} \qquad (3)\]

股价延迟定义为： \[\text{PriceDelay}_{i,t} = 1 - \frac{R^2_{\text{restricted}}}{R^2_{\text{unrestricted}}} \qquad (4)\]

取值范围为{[}0,1{]}，值越大表示股价对市场信息的反应越滞后，定价效率越低。若H1成立，DU应与PriceDelay呈负向关联。

作为辅助指标，本文同时采用股价同步性SYNCH = ln(R²/(1-R²))，其中R²来自个股收益率对市场收益率和行业收益率的回归。PriceDelay衡量价格对市场信息的吸收速度，SYNCH反映市场与行业共同因素对个股收益率的解释比例。若数据要素利用改善了信息环境，SYNCH上升可能反映系统性信息定价更快或特质噪声下降，但其经济含义不同于PriceDelay，因此本文仅将其作为辅助结果报告，而不将其作为H1的镜像验证。

自变量为数据要素利用。主指标DU\_kw和语义增强指标DU\_llm的定义和构造方法已在上一节详述；DUevent、DUchain\_count、DUclosedloop、DUcore和DUkw\_mda则作为支持性增强测度，在后文用于检验``更完整、更嵌入''的数据利用表述是否具有与基准指标一致的定价含量。

控制变量共11个，包括企业规模、财务状况和治理特征：

\textbf{表2 控制变量定义}

{\def\LTcaptype{} % do not increment counter
\begin{longtable}[]{@{}ll@{}}
\toprule\noalign{}
变量 & 定义 \\
\midrule\noalign{}
\endhead
\bottomrule\noalign{}
\endlastfoot
Size & 总市值的自然对数 \\
Lev & 总负债/总资产 \\
ROA & 净利润/总资产 \\
TobinQ & (股权市值+负债账面值)/总资产 \\
Age & ln(当年-上市年份+1) \\
Growth & 营收增长率 \\
IndepRatio & 独立董事人数/董事会总人数 \\
Dual & 董事长兼任CEO取1，否则为0 \\
Top1Share & 第一大股东持股比例 \\
SOE & 国有企业取1，否则为0 \\
CFO & 经营活动现金流/总资产 \\
\end{longtable}
}

控制变量参照股价延迟研究的通行做法选取。由于机制分析关注的是数据要素利用如何影响分析师预测分歧、现金流波动和供应链集中度，这些变量本身不纳入基准控制，以避免``坏控制''问题（Angrist and Pischke, 2009）。

标准误在行业×年份层面聚类，以校正同行业企业在同一年份受共同技术扩散冲击导致的截面相关。稳健性检验中报告了企业和年份双向聚类的结果，结论不变。

\subsection{（四）模型设定}\label{ux56dbux6a21ux578bux8bbeux5b9a}

为检验数据要素利用对资本市场定价效率的影响，构建如下基准回归模型：

\[PriceDelay_{i,t} = \alpha_0 + \alpha_1 DU_{i,t} + \alpha_2 Control_{i,t} + \mu_i + \lambda_t + \varepsilon_{i,t} \qquad (5)\]

其中，PriceDelay为企业\(i\)在第\(t\)年的股价延迟；DU为数据要素利用指标，分别以DU\_kw和DU\_llm进入回归；Control为11个控制变量向量；\(\mu_i\)和\(\lambda_t\)分别为企业固定效应和年份固定效应；\(\varepsilon_{i,t}\)为随机扰动项。标准误在行业×年份层面聚类。若H1成立，\(\alpha_1\)应显著为负。

内生性方面，本文采用两个工具变量：IV1为同行业跨省企业DU的留一均值，利用行业层面技术扩散的外生变异；IV2为省级大数据发展总指数（2016年截面）与年份趋势的交互项，利用地理维度基础设施供给的外生变异。排除限制的论证和诊断检验在第四节报告。Oster界检验、Heckman两阶段法和安慰剂检验作为辅助识别策略。

传导机制的检验模型为：

\[Channel_{i,t} = \beta_0 + \beta_1 DU_{i,t} + \beta_2 Control_{i,t} + \mu_i + \lambda_t + \varepsilon_{i,t} \qquad (6)\]

其中Channel为渠道变量，定义如下：

\textbf{表3 渠道变量定义}

{\def\LTcaptype{} % do not increment counter
\begin{longtable}[]{@{}ll@{}}
\toprule\noalign{}
渠道变量 & 定义 \\
\midrule\noalign{}
\endhead
\bottomrule\noalign{}
\endlastfoot
ForecastDisp 分析师预测分歧 & 同一企业年度内至少3条盈利预测的标准差 \\
CashFlowVol 现金流波动率 & 经营活动现金流/总资产的三年滚动标准差 \\
SCConc 综合供应链集中度 & 前五大客户销售占比与前五大供应商采购占比的均值 \\
\end{longtable}
}

遵循江艇（2022）的建议，采用路径系数的直接呈现来识别传导渠道的存在性，同时使用DU\_kw和DU\_llm两套测度进行交叉验证。双测度下方向一致且均显著的渠道视为更稳健的路径证据；仅在一套测度下显著的渠道一并报告，但对其解释保持审慎。渠道变量均经1\%/99\%缩尾处理。

异质性分析方面，主文在统一滞后规格下优先报告三条仍具有较清晰解释力的边界条件：政策阶段差异、高分析师关注与非主板企业。具体地，政策阶段以2022年为切点；高分析师关注按分析师覆盖的全样本中位数分组；MainBoard取1表示主板公司、0表示非主板公司。数字经济核心产业、替代切点以及其他补充分组结果仍放在附录中，作为补充证据报告。考虑到统一滞后规格后部分产业差异不再稳健，本文对异质性结论主要采用交互项检验；政策阶段结果另辅以费舍尔组合检验并使用更审慎的解释口径。

\section{四、实证检验}\label{ux56dbux5b9eux8bc1ux68c0ux9a8c}

\subsection{（一）描述性统计}\label{ux4e00ux63cfux8ff0ux6027ux7edfux8ba1}

表4报告了主要变量的描述性统计结果（见表4）。样本包含43,735个企业-年份观测。股价延迟均值为0.112，中位数为0.068，多数企业的信息融入效率较高。DU\_kw均值为1.219，标准差为1.861，企业间差异较大。DU\_llm均值为1.321，标准差为1.119。其余控制变量的分布均在合理范围内。

\textbf{表4 变量的描述性统计}

{\def\LTcaptype{} % do not increment counter
\begin{longtable}[]{@{}lllllll@{}}
\toprule\noalign{}
变量 & N & 均值 & 标准差 & 最小值 & 中位数 & 最大值 \\
\midrule\noalign{}
\endhead
\bottomrule\noalign{}
\endlastfoot
PriceDelay & 43,735 & 0.112 & 0.123 & 0.005 & 0.068 & 0.665 \\
SYNCH & 43,677 & -0.505 & 0.834 & -2.954 & -0.449 & 1.229 \\
DU\_kw & 43,735 & 1.219 & 1.861 & 0.000 & 0.572 & 11.064 \\
DU\_llm & 43,735 & 1.321 & 1.119 & 0.000 & 1.086 & 6.799 \\
Size & 43,735 & 22.597 & 0.971 & 20.948 & 22.435 & 25.608 \\
Lev & 43,735 & 0.425 & 0.207 & 0.055 & 0.417 & 0.916 \\
ROA & 43,735 & 0.030 & 0.067 & -0.280 & 0.033 & 0.191 \\
TobinQ & 43,735 & 2.005 & 1.283 & 0.829 & 1.590 & 8.451 \\
Age & 43,735 & 2.159 & 0.733 & 0.693 & 2.303 & 3.219 \\
Growth & 43,735 & 0.134 & 0.369 & -0.591 & 0.083 & 2.157 \\
IndepRatio & 43,735 & 0.378 & 0.054 & 0.333 & 0.364 & 0.571 \\
Dual & 43,735 & 0.297 & 0.457 & 0.000 & 0.000 & 1.000 \\
Top1Share & 43,735 & 33.420 & 14.822 & 8.126 & 31.000 & 74.000 \\
SOE & 43,735 & 0.329 & 0.470 & 0.000 & 0.000 & 1.000 \\
CFO & 43,735 & 0.046 & 0.068 & -0.161 & 0.045 & 0.239 \\
\end{longtable}
}

LLM语义评分中，得分2及以上（实质利用）的企业-年份观测占46.8\%，得分3（深度利用）仅占5.4\%。

\subsection{（二）基准回归结果}\label{ux4e8cux57faux51c6ux56deux5f52ux7ed3ux679c}

表5报告了数据要素利用对股价延迟的OLS基准回归结果（见表5）。模型(1)(2)使用基准处理变量DU\_kw，模型(3)(4)使用语义增强指标DU\_llm，模型(5)使用LLM二值化指标。所有模型均控制企业和年份固定效应，聚类标准误设在行业×年份层面。

\textbf{表5 基准回归检验}

{\def\LTcaptype{} % do not increment counter
\begin{longtable}[]{@{}llllll@{}}
\toprule\noalign{}
& (1) & (2) & (3) & (4) & (5) \\
\midrule\noalign{}
\endhead
\bottomrule\noalign{}
\endlastfoot
变量 & PriceDelay & PriceDelay & PriceDelay & PriceDelay & PriceDelay \\
DU\_kw & -0.0052*** & -0.0046*** & & & \\
& (0.0010) & (0.0009) & & & \\
DU\_llm & & & -0.0037*** & -0.0036*** & \\
& & & (0.0008) & (0.0008) & \\
llm\_binary & & & & & -0.0033*** \\
& & & & & (0.0013) \\
Size & & 0.0063*** & & 0.0060** & 0.0055** \\
& & (0.0024) & & (0.0024) & (0.0024) \\
Lev & & 0.0225*** & & 0.0228*** & 0.0225*** \\
& & (0.0063) & & (0.0063) & (0.0063) \\
ROA & & -0.0634*** & & -0.0612*** & -0.0614*** \\
& & (0.0147) & & (0.0147) & (0.0147) \\
TobinQ & & 0.0116*** & & 0.0117*** & 0.0119*** \\
& & (0.0011) & & (0.0011) & (0.0011) \\
Age & & -0.0331*** & & -0.0346*** & -0.0348*** \\
& & (0.0038) & & (0.0039) & (0.0039) \\
Growth & & 0.0129*** & & 0.0130*** & 0.0130*** \\
& & (0.0020) & & (0.0020) & (0.0020) \\
IndepRatio & & -0.0159 & & -0.0149 & -0.0142 \\
& & (0.0155) & & (0.0155) & (0.0155) \\
Dual & & -0.0021 & & -0.0021 & -0.0021 \\
& & (0.0016) & & (0.0016) & (0.0016) \\
Top1Share & & 0.0000 & & 0.0000 & 0.0000 \\
& & (0.0001) & & (0.0001) & (0.0001) \\
SOE & & 0.0051 & & 0.0052 & 0.0053 \\
& & (0.0040) & & (0.0040) & (0.0040) \\
CFO & & 0.0047 & & 0.0051 & 0.0055 \\
& & (0.0111) & & (0.0111) & (0.0111) \\
Firm FE & YES & YES & YES & YES & YES \\
Year FE & YES & YES & YES & YES & YES \\
N & 43,721 & 43,721 & 43,721 & 43,721 & 43,721 \\
R² & 0.409 & 0.419 & 0.408 & 0.418 & 0.418 \\
\end{longtable}
}

以加入全部控制变量的模型(2)为基准，DU\_kw每增加一个标准差1.861，股价延迟预期下降约7.7\%。模型(4)中DU\_llm的系数为-0.0036，经济意义同样显著。模型(5)中llm\_binary的系数为-0.0033（t=-2.61），表明即便将语义深度简化为二分类变量，其与股价延迟之间的负向关系仍然显著。两套测度在PriceDelay上的结果方向一致，其中DU\_kw作为基准处理变量承担主识别角色，DU\_llm提供语义层面的并行验证。

以SYNCH为辅助指标时，DU\_kw系数为+0.0189（t=3.88），DU\_llm系数为+0.0176（t=3.98），均显著为正，说明收益率中由市场和行业共同因素解释的比例提高。由于SYNCH的经济含义不同于PriceDelay，本文仅将其视为辅助结果，而不据此单独验证H1。

\subsection{（三）构念边界与测度补强}\label{ux4e09ux6784ux5ff5ux8fb9ux754cux4e0eux6d4bux5ea6ux8865ux5f3a}

构念边界与测度补强检验的核心问题是：本文捕捉到的究竟是更贴近数据利用的信息，还是更宽泛的数字叙事强度。为此，依次采用篇幅控制、去泛化词口径、长度标准化语义深度和GenericNarr联合回归进行检验。

\textbf{表6 构念边界与测度补强}

\begin{landscape}
\begingroup
\setlength{\tabcolsep}{3pt}
\scriptsize
{\def\LTcaptype{} % do not increment counter
\begin{longtable}[]{@{}>{\raggedright\arraybackslash}p{3.0cm}*{7}{>{\centering\arraybackslash}p{2.0cm}}@{}}
\toprule\noalign{}
& (1) & (2) & (3) & (4) & (5) & (6) & (7) \\
\midrule\noalign{}
\endhead
\bottomrule\noalign{}
\endlastfoot
模型 & DUkw+Total & DUkw+MDA & DUllm+Total & Lenstd+Total & DUkw+Narr & ResidDU & StrictDU \\
因变量 & PriceDelay & PriceDelay & PriceDelay & PriceDelay & PriceDelay & PriceDelay & PriceDelay \\
DU\_kw\_lag & -0.0038*** & -0.0039*** & & & -0.0027 & & \\
& (0.0010) & (0.0010) & & & (0.0021) & & \\
DU\_llm\_lag & & & -0.0042*** & & & & \\
& & & (0.0009) & & & & \\
DU\_llm\_lenstd\_lag & & & & -0.0184*** & & & \\
& & & & (0.0032) & & & \\
ln\_total\_chars\_lag & -0.0171*** & & -0.0150** & -0.0165*** & & & \\
& (0.0060) & & (0.0060) & (0.0059) & & & \\
ln\_mda\_chars\_lag & & 0.0003 & & & & & \\
& & (0.0003) & & & & & \\
GenericNarr\_lag & & & & & -0.0017 & & \\
& & & & & (0.0023) & & \\
DU\_kw\_resid\_lag & & & & & & -0.0027 & \\
& & & & & & (0.0021) & \\
DU\_kw\_strict\_lag & & & & & & & -0.0077*** \\
& & & & & & & (0.0029) \\
控制变量 & YES & YES & YES & YES & YES & YES & YES \\
Firm FE & YES & YES & YES & YES & YES & YES & YES \\
Year FE & YES & YES & YES & YES & YES & YES & YES \\
N & 37,294 & 37,294 & 37,294 & 37,294 & 37,294 & 37,294 & 37,294 \\
\end{longtable}
}
\endgroup
\end{landscape}

表6显示，首先，在控制年报总字数和MD\&A字数后，DU\_kw\_lag的系数分别为-0.0038（t=-3.80）和-0.0039（t=-3.89），仍在1\%水平上显著为负，说明主结果并非单纯由披露篇幅更长的企业驱动。其次，DU\_llm\_lag在控制总字数后仍显著为负，长度标准化后的DU\_llm\_lenstd\_lag系数进一步扩大至-0.0184（t=-5.80），表明语义深度在控制篇幅污染后仍保留定价信息。

但更严格的构念边界检验并未支持``数据要素利用''与一般数字叙事完全分离的强主张。在DU\_kw\_lag与GenericNarr\_lag同时进入时，两者系数均不显著；FE口径下的VIF约为7.6，说明二者存在较强重叠。进一步将DU\_kw\_lag中可被GenericNarr解释的部分剔除后，残差项DU\_kw\_resid\_lag不再显著。这意味着本文更稳妥的表述应是：当前指标较一般数字叙事更贴近``数据资源实际利用''，但并不能被理解为与更宽泛数字化叙事完全正交的纯净构念。

与此同时，去泛化词后的DU\_kw\_strict\_lag系数为-0.0077（t=-2.61），仍在1\%水平上显著为负，说明核心关系并不完全依赖``数字化转型、人工智能、大数据''等高频泛化词。综合篇幅控制、长度标准化语义深度与去泛化词结果，本文认为主结论具备较强的测度稳健性；但对其概念独立性的表述仍需保持审慎。

仅有边界检验还不足以说明测度是否真正朝``数据要素利用''推进。本文进一步利用现有五维词典和章节特征构造增强口径，包括规则识别的实质利用指标DUevent、链条完整度指标DUchain\_count、闭环式表述指标DUclosedloop、窄口径核心强度指标DUcore以及MD\&A章节限定指标DUkw\_mda。表6A报告了这些增强口径在统一滞后规格下的结果。

\textbf{表6A 增强测度与章节限定口径检验}

\begin{landscape}
\begingroup
\setlength{\tabcolsep}{3pt}
\scriptsize
{\def\LTcaptype{} % do not increment counter
\begin{longtable}[]{@{}>{\raggedright\arraybackslash}p{3.2cm}*{6}{>{\centering\arraybackslash}p{2.3cm}}@{}}
\toprule\noalign{}
& (1) & (2) & (3) & (4) & (5) & (6) \\
\midrule\noalign{}
\endhead
\bottomrule\noalign{}
\endlastfoot
模型 & DUevent & DUchain & DUclosed & DUcore & DUkw\_mda & DUclosed+Narr \\
因变量 & PriceDelay & PriceDelay & PriceDelay & PriceDelay & PriceDelay & PriceDelay \\
DUevent\_lag & -0.0023 & & & & & \\
& (0.0016) & & & & & \\
DUchain\_count\_lag & & -0.0018** & & & & \\
& & (0.0009) & & & & \\
DUclosedloop\_lag & & & -0.0083*** & & & -0.0051* \\
& & & (0.0030) & & & (0.0029) \\
DUcore\_lag & & & & -0.0095*** & & \\
& & & & (0.0027) & & \\
DUkw\_mda\_lag & & & & & -0.0006*** & \\
& & & & & (0.0002) & \\
GenericNarr\_lag & & & & & & -0.0046*** \\
& & & & & & (0.0012) \\
控制变量 & YES & YES & YES & YES & YES & YES \\
Firm FE & YES & YES & YES & YES & YES & YES \\
Year FE & YES & YES & YES & YES & YES & YES \\
N & 37,294 & 37,294 & 37,294 & 37,294 & 37,294 & 37,294 \\
\end{longtable}
}
\endgroup
\end{landscape}

表6A显示，增强测度中并非所有口径都能提供同等强度的支持。首先，规则识别的DUevent\_lag系数为-0.0023（t=-1.51），方向为负但统计上不显著，说明在当前归档特征表条件下，单纯依赖rule-based实质利用计数尚不足以替代基准广度指标。其次，DUchain\_count\_lag的系数为-0.0018（t=-2.13），DUclosedloop\_lag的系数为-0.0083（t=-2.74），均显著为负，说明越完整、越接近``数据资源---加工应用---治理/价值实现''闭环的数据利用表述，与更快的价格发现相关。再次，DUcore\_lag和DUkw\_mda\_lag也分别显著为负，表明更窄的经营/治理口径以及经营讨论章节中的数据利用披露同样具有定价含量。

但这些增强测度并未形成足以替代DU\_kw的新主指标。附加联合回归结果表明，DUchain\_count、DUcore和DUkw\_mda在与DU\_kw或GenericNarr同时进入后大多不再稳健；相较之下，DUclosedloop在与GenericNarr联合回归时仍保持-0.0051的负系数（t=-1.77），但仅达到10\%水平。基于这一证据，本文对构念升级采取更克制的写法：本文不能宣称已经从根本上摆脱了宽泛数字叙事的共线性问题，但可以更有把握地说明，资本市场对``更完整、更嵌入''的数据利用表述确实给出了方向一致的积极反应。

\subsection{（四）时间顺序与内生性识别}\label{ux56dbux65f6ux95f4ux987aux5e8fux4e0eux5185ux751fux6027ux8bc6ux522b}

基准回归中的内生性风险主要来自两个方面。其一，信息环境较好的企业可能同时呈现高数据利用和低股价延迟，构成遗漏变量偏差。其二，定价效率较高的企业可能更有动力披露数据利用信息，形成反向因果。为增强时间顺序的清晰性，本文首先报告滞后规格，并在此基础上结合工具变量、Oster界检验、Heckman两阶段法和安慰剂检验进行识别。

\textbf{表7 时间顺序与内生性识别}

\begin{landscape}
\begingroup
\setlength{\tabcolsep}{3pt}
\scriptsize
{\def\LTcaptype{} % do not increment counter
\begin{longtable}[]{@{}>{\raggedright\arraybackslash}p{3.0cm}*{5}{>{\centering\arraybackslash}p{2.55cm}}@{}}
\toprule\noalign{}
& (1) & (2) & (3) & (4) & (5) \\
\midrule\noalign{}
\endhead
\bottomrule\noalign{}
\endlastfoot
模型 & OLS & Lag OLS & Lag IV & IV1:Peer & IV2:Dig×Year \\
DU\_kw & -0.0046*** & -0.0039*** & -0.0163*** & -0.0143*** & -0.0352*** \\
& (0.0009) & (0.0010) & (0.0053) & (0.0047) & (0.0114) \\
First-stage F & & & 177.8 & 251.3 & 72.7 \\
控制变量 & YES & YES & YES & YES & YES \\
Firm FE & YES & YES & YES & YES & YES \\
Year FE & YES & YES & YES & YES & YES \\
N & 43,695 & 37,278 & 37,198 & 43,612 & 43,695 \\
\end{longtable}
}
\endgroup
\end{landscape}

表7显示，以滞后一期的DU\_kw解释当期PriceDelay后，OLS系数为-0.0039（t=-3.90），方向和显著性与同期规格一致；使用滞后跨省peer作为IV时，系数为-0.0163（t=-3.10），一阶段F=177.8，说明在更清晰的时间顺序下，数据要素利用与更低股价延迟之间的关系依然成立。同期两套工具变量的估计也均显著为负，且一阶段统计量远超弱工具变量常用临界值，表明行业技术扩散与地理基础设施供给两类外生变异均支持基准结论。Durbin-Wu-Hausman检验在三个IV规格下均拒绝外生性假设（p值分别为0.015、0.005和0.006），意味着将内生性问题纳入考虑是必要的。

采用Oster（2019）提出的δ界检验方法评估遗漏变量的潜在影响。以Rmax=1.3R²为基准，δ*=96.3；以更保守的Rmax=2R²-R²\_short为基准，δ*=7.5。两个δ*均远超临界值1，意味着若要使DU\_kw的系数降为零，遗漏变量对因变量的解释能力需要达到已纳入控制变量的96倍，或在更保守口径下达到7.5倍，这在经验上不太可能。

Heckman两阶段法进一步检验了大量零值可能带来的样本选择偏差。约23\%的企业-年份观测中DU\_kw=0。第一阶段以Probit模型估计企业年报出现数据要素相关表述的概率，第二阶段在DU\_kw\textgreater0的子样本中加入逆米尔斯比IMR控制选择偏差。结果显示，DU\_kw的系数为-0.0040（t=-4.39），与基准回归几乎一致；IMR的系数为-0.123（t=-6.23），在1\%水平上显著，说明样本选择偏差存在，但控制之后核心结论不变。

安慰剂检验则从反向因果角度进行检验。如果数据要素利用对股价延迟的影响确实不是伪回归，那么\(t+1\)期的DU\_kw不应影响\(t\)期的PriceDelay。将DU\_\{i,t\}和DU\_\{i,t+1\}同时纳入回归后，当期系数为-0.0050（t=-4.64），未来期系数为0.0001（t=0.11），完全不显著。未来期DU对当期定价效率没有解释力，与反向因果的预测不符。综合滞后规格、工具变量、Oster界检验、Heckman修正和安慰剂检验，本文认为多种识别练习共同支持数据要素利用降低股价延迟的因果解释。

\subsection{（五）双重机器学习估计}\label{ux4e94ux53ccux91cdux673aux5668ux5b66ux4e60ux4f30ux8ba1}

为放松线性函数形式假定，采用Chernozhukov等（2018）的双重机器学习（DML）框架重新估计处理效应，以LightGBM为第一阶段学习器，Mundlak（1978）均值替代法处理企业固定效应，控制变量扩展至27个。

\textbf{表8 DML-PLR回归结果}

{\def\LTcaptype{} % do not increment counter
\begin{longtable}[]{@{}lllllll@{}}
\toprule\noalign{}
& Y=PriceDelay & & & Y=SYNCH & & \\
\midrule\noalign{}
\endhead
\bottomrule\noalign{}
\endlastfoot
处理变量D & 系数 & 标准误 & t值 & 系数 & 标准误 & t值 \\
DU\_kw & -0.0029*** & (0.0004) & {[}-6.50{]} & +0.0277*** & (0.0029) & {[}9.65{]} \\
DU\_kw\_ln & -0.0040*** & (0.0006) & {[}-7.19{]} & & & \\
DU\_sub\_ln & -0.0039*** & (0.0005) & {[}-7.21{]} & & & \\
DU\_llm & -0.0042*** & (0.0007) & {[}-6.35{]} & +0.0326*** & (0.0041) & {[}7.98{]} \\
llm\_binary & -0.0042*** & (0.0012) & {[}-3.46{]} & +0.0435*** & (0.0075) & {[}5.80{]} \\
控制变量+CRE & YES & & & YES & & \\
N & 35,666 & & & 35,614 & & \\
\end{longtable}
}

表8显示，DU\_kw的DML系数为-0.0029（t=-6.50），DU\_llm的系数为-0.0042（t=-6.35），均在1\%水平上显著为负，说明在放松线性函数形式并引入更高维控制后，核心结果依然稳健。SYNCH列的结果同样为正，但仍仅作为辅助信息报告，而不承担主结论验证职责。DML样本约35,700，因27个控制变量对完整观测的要求高于OLS样本，样本量有所下降，但估计方向与基准回归保持一致。

\subsection{（六）稳健性检验}\label{ux516dux7a33ux5065ux6027ux68c0ux9a8c}

表9报告了多维度的稳健性检验结果（见表9）。

（1）替换固定效应。在保留企业固定效应的基础上，以行业×年份联合固定效应替代单独年份固定效应，以吸收同一行业在特定年份面临的共同冲击，DU\_kw系数为-0.0029（t=-4.91），结论不变。

（2）调整样本范围。剔除末期年份（2024年）后系数为-0.0052（t=-5.43），排除了样本末端效应的干扰。剔除证监会行业分类中的信息技术业后系数为-0.0046（t=-5.29），排除了结论由数据密集型行业驱动的可能。

（3）滞后控制变量。以\(t-1\)期控制变量替代同期变量，缓解控制变量与处理变量的同期内生性，系数为-0.0040（t=-4.38）。

（4）倾向得分匹配（Rosenbaum和Rubin，1983）。以DU\_kw高于行业-年份中位数为处理组，Logit模型估计倾向得分，采用最近邻1:1匹配（卡尺0.05）。图1展示了匹配前后处理组与对照组的倾向得分核密度分布，匹配后两组分布趋于一致，协变量平衡良好。匹配样本上系数为-0.0039（t=-4.61），与基准估计方向和大小一致。

\textbf{图1 匹配前后核密度曲线}

\begin{figure}[H]
\centering
\includegraphics[width=0.78\linewidth]{results/v16_tables/psm_density.png}
\end{figure}

（5）双向聚类标准误（Cameron等，2011）。将聚类层级从行业×年份扩展至企业和年份双向聚类，以应对企业层面和时间层面的双重相关性，系数为-0.0046（t=-3.81），标准误增大但结论不变。

（6）替换自变量。以实质利用指标DU\_sub\_ln替代DU\_kw，系数为-0.0042（t=-5.49），表明结论不依赖于特定的变量构造方式。

\textbf{表9 稳健性检验}

\begin{landscape}
\begingroup
\setlength{\tabcolsep}{3pt}
\scriptsize
{\def\LTcaptype{} % do not increment counter
\begin{longtable}[]{@{}>{\raggedright\arraybackslash}p{3.0cm}*{7}{>{\centering\arraybackslash}p{2.05cm}}@{}}
\toprule\noalign{}
& (1)FE & (2)尾年 & (3)IT & (4)滞后 & (5)PSM & (6)双聚类 & (7)替换 \\
\midrule\noalign{}
\endhead
\bottomrule\noalign{}
\endlastfoot
& 行业×年FE & 剔除2024 & 剔除IT & 滞后控制 & 倾向匹配 & 双向聚类 & DU\_sub\_ln \\
变量 & PriceDelay & PriceDelay & PriceDelay & PriceDelay & PriceDelay & PriceDelay & PriceDelay \\
DU\_kw & -0.0029*** & -0.0052*** & -0.0046*** & -0.0040*** & -0.0039*** & -0.0046*** & \\
(DU\_sub\_ln) & & & & & & & -0.0042*** \\
& (0.0006) & (0.0010) & (0.0009) & (0.0009) & (0.0008) & (0.0012) & (0.0008) \\
控制变量 & YES & YES & YES & YES & YES & YES & YES \\
Firm FE & YES & YES & YES & YES & YES & YES & YES \\
Year FE & NO & YES & YES & YES & YES & YES & YES \\
Ind×Year FE & YES & NO & NO & NO & NO & NO & NO \\
N & 43,656 & 38,812 & 40,613 & 37,294 & 35,929 & 43,721 & 43,721 \\
\end{longtable}
}
\endgroup
\end{landscape}

此外，以SYNCH替代PriceDelay时，DU\_kw的OLS系数为+0.0189（t=3.88），与基准回归中的辅助指标结果一致，即数据要素利用提高了收益率中由市场和行业共同因素解释的比例。由于该指标的经济含义不同于PriceDelay，本文对其解释保持审慎。总体而言，上述稳健性检验与前文识别结果共同表明，数据要素利用与更低股价延迟之间的负向关系并不依赖于特定模型设定或样本处理方式。

\section{五、进一步分析}\label{ux4e94ux8fdbux4e00ux6b65ux5206ux6790}

\subsection{（一）传导机制检验}\label{ux4e00ux4f20ux5bfcux673aux5236ux68c0ux9a8c}

借鉴江艇（2022）的建议，以渠道变量为因变量、数据要素利用及控制变量为自变量进行回归，检验信息从年报向股价传导的具体路径。三个渠道变量分别为分析师预测分歧（ForecastDisp，同一企业年度内至少3条盈利预测的标准差）、现金流波动率（CashFlowVol，经营性现金流与总资产之比的三年滚动标准差）和综合供应链集中度（SCConc，前五大客户销售占比与前五大供应商采购占比的均值）。所有渠道变量经1\%/99\%缩尾处理，模型控制企业与年份固定效应和11个控制变量，标准误在行业×年份水平聚类。

\textbf{1. 分析师预测分歧渠道}

数据要素利用涉及技术架构、业务流程和价值实现方式，投资者对其经济含义的理解存在天然门槛。当企业在年报中提供更丰富的数据利用信息时，分析师更容易围绕同一组经营事实更新预期，预测分歧因而可能下降。相较于分析师覆盖的``数量''，预测分歧衡量的是分析师群体对同一企业判断的``共识程度''，更接近信息解释是否趋于一致这一问题。

表10第(1)(2)列结果显示，DU\_kw对分析师预测分歧（ForecastDisp）的系数为-0.0065（t=-4.08），在1\%水平上显著。数据要素利用程度越高，分析师之间的盈利预测分歧越小。DU\_llm的系数不显著，表明关键词频率所承载的广度信息更有助于缩小分析师间的认知差异，而语义深度的评估尚未被分析师群体稳定吸收。考虑到该渠道仅在DU\_kw下显著，且样本量仅为22,685，该结果更适合作为补充性的信息质量证据。

\textbf{2. 经营风险渠道}

数据驱动的运营管理有助于提高需求预测精度、优化库存配置和降低生产波动，从而平滑企业经营层面的现金流。经营层面的不确定性越低，企业基本面信号越稳定，投资者对未来现金流的估值也越容易趋于收敛，信息向价格的传导因而更加顺畅。

表10第(3)(4)列结果显示，DU\_kw对现金流波动率（CashFlowVol）的系数为-0.0005（t=-2.45），DU\_llm的系数为-0.0005（t=-2.05），均在5\%水平上显著。数据要素利用降低了企业经营层面的现金流波动，双测度方向一致且均显著，经营风险渠道得到较为稳健的支持。

\textbf{3. 供应链集中度渠道}

供应链过度集中意味着企业对少数交易对手的依赖较高，一旦关键客户或供应商发生变化，经营基本面将面临更大的冲击，市场对企业价值评估的不确定性也会随之上升。数据要素利用可以提升企业对订单、库存和交易伙伴的识别与匹配能力，促进客户和供应商结构的动态优化，使供应链关系更为分散和稳定。

表10第(5)(6)列结果显示，DU\_kw对综合供应链集中度（SCConc）的系数为-0.277（t=-4.00），DU\_llm的系数为-0.460（t=-6.77），均在1\%水平上高度显著。数据要素利用显著降低了供应链集中度，双测度均高度显著，供应链渠道成立。

\textbf{表10 传导机制检验}

{\def\LTcaptype{} % do not increment counter
\begin{longtable}[]{@{}lllllll@{}}
\toprule\noalign{}
& (1) & (2) & (3) & (4) & (5) & (6) \\
\midrule\noalign{}
\endhead
\bottomrule\noalign{}
\endlastfoot
变量 & ForecastDisp & ForecastDisp & CashFlowVol & CashFlowVol & SCConc & SCConc \\
DU\_kw & -0.0065*** & & -0.0005** & & -0.277*** & \\
& (0.0016) & & (0.0002) & & (0.069) & \\
DU\_llm & & +0.0010 & & -0.0005** & & -0.460*** \\
& & (0.0017) & & (0.0002) & & (0.068) \\
控制变量 & YES & YES & YES & YES & YES & YES \\
Firm FE & YES & YES & YES & YES & YES & YES \\
Year FE & YES & YES & YES & YES & YES & YES \\
N & 22,685 & 22,685 & 43,721 & 43,721 & 42,332 & 42,332 \\
\end{longtable}
}

综合来看，现金流波动率和供应链集中度是双测度下更稳健的传导路径证据，可以视为本文机制识别的主体结果；分析师预测分歧则提供了补充性的信息质量证据。三条路径共同表明，数据要素利用不仅改变了资本市场对信息的解释方式，也与企业经营稳定性和供应链韧性的改善相联系。

\subsection{（二）广度---深度错配与概念性叙事}\label{ux4e8cux5e7fux5ea6ux6df1ux5ea6ux9519ux914dux4e0eux6982ux5ff5ux6027ux53d9ux4e8b}

前文的构念边界与增强测度检验显示，DU\_kw与更宽泛数字叙事存在交叠，但链条完整度、闭环式表述和长度标准化后的语义深度均与更快的价格发现方向一致。进一步的问题是：资本市场究竟奖励``提及得多''的概念性叙事，还是奖励``说得具体、用得更深''的实质利用披露。为此，本文构造WashGap = z(DU\_kw) - z(DU\_llm\_lenstd)，衡量关键词广度相对于语义深度的超额部分；同时构造BroadShallow虚拟变量，将当年DU\_kw位于上三分位而DU\_llm\_lenstd位于下三分位的企业识别为``提得多但说得浅''的样本。所有检验均采用与主文一致的滞后规格。

\textbf{表11 广度---深度错配与概念性叙事}

{\def\LTcaptype{} % do not increment counter
\begin{longtable}[]{@{}llll@{}}
\toprule\noalign{}
& (1) & (2) & (3) \\
\midrule\noalign{}
\endhead
\bottomrule\noalign{}
\endlastfoot
模型 & WashGap & BroadShallow & 联合回归 \\
因变量 & PriceDelay & PriceDelay & PriceDelay \\
WashGap\_lag & +0.0035** & & \\
& (0.0017) & & \\
BroadShallow\_lag & & +0.0046 & \\
& & (0.0036) & \\
DU\_kw\_lag & & & -0.0016 \\
& & & (0.0012) \\
DU\_llm\_lenstd\_lag & & & -0.0146*** \\
& & & (0.0039) \\
控制变量 & YES & YES & YES \\
Firm FE & YES & YES & YES \\
Year FE & YES & YES & YES \\
N & 37,294 & 37,294 & 37,294 \\
\end{longtable}
}

表11第(1)列显示，WashGap\_lag的系数为+0.0035（t=2.09），在5\%水平上显著为正，说明当企业在关键词广度上``说得更多''，但语义深度并未同步提高时，股价对相关信息的吸收反而更慢。也就是说，概念性叙事与实质利用之间的错配会加剧价格发现摩擦。

第(2)列中，BroadShallow\_lag的系数为+0.0046（t=1.28），方向与预期一致但统计上不显著。这意味着以分位数离散化构造的``宽而浅''虚拟变量可以提供方向一致的补充证据，但其识别强度弱于连续型错配指标WashGap。

更关键的是，第(3)列同时纳入DU\_kw\_lag与DU\_llm\_lenstd\_lag后，DU\_llm\_lenstd\_lag的系数为-0.0146（t=-3.76），在1\%水平上显著，而DU\_kw\_lag不再显著。该结果表明，在控制关键词广度之后，语义深度仍保留额外的负向解释力。结合WashGap的正向结果，以及前文DUclosedloop和DUchain\_count的负向证据，可以将这一扩展分析概括为：资本市场更倾向于奖励具有实质内容、经营嵌入和链条完整性的数据利用披露，而非单纯更频繁但较浅层的概念性提及。由于这一部分属于扩展分析，本文据此补充主文叙事，但不以其替代前文基准回归和识别结果。

\section{六、异质性分析}\label{ux516dux5f02ux8d28ux6027ux5206ux6790}

为考察数据要素利用的定价效率改善效应在不同制度环境下的差异，本文在统一滞后规格下报告三条仍具有较清晰解释力的异质性结果：政策阶段，以及高分析师关注和上市板块两条补充性边界条件。所有模型均控制企业与年份固定效应和11个控制变量，标准误在行业×年份水平聚类。政策阶段结果同时结合交互项检验和费舍尔组合检验，高分析师关注与上市板块则以交互项检验为主。

\subsection{（一）政策阶段效应}\label{ux4e00ux653fux7b56ux9636ux6bb5ux6548ux5e94}

相较于2020或2021，2022年``数据二十条''更接近数据要素制度密集推进的起点，也更适合作为主文中的政策阶段切点。其背后的逻辑并非机械地将估计结果解释为某一单一政策的即时效应，而是将2022年之后理解为``数据利用披露逐步制度化、市场预期更充分''的阶段。若数据利用信息在较早阶段更具新颖性和辨识度，则其对价格发现的边际贡献应在2022年前更强。

表12第(1)(2)列结果显示，2011至2021年样本中DU\_kw系数为-0.0066（t=-5.13），2022至2024年样本中系数为-0.0018（t=-0.85，不显著），交互项检验P值为0.031，费舍尔组合检验得到相同方向的结论。这意味着数据要素利用降低股价延迟的效应主要集中在制度化推进之前的阶段：在相关披露尚未成为更普遍表达时，年报中的数据利用信息更容易构成有效的信息更新；而当市场对该类叙事已有更高先验预期后，其边际定价价值有所减弱。

\textbf{表12 政策阶段异质性检验}

{\def\LTcaptype{} % do not increment counter
\begin{longtable}[]{@{}lll@{}}
\toprule\noalign{}
& (1) & (2) \\
\midrule\noalign{}
\endhead
\bottomrule\noalign{}
\endlastfoot
变量 & 2011-2021 & 2022-2024 \\
& PriceDelay & PriceDelay \\
DU\_kw & -0.0066*** & -0.0018 \\
& (0.0013) & (0.0021) \\
控制变量 & YES & YES \\
Firm FE & YES & YES \\
Year FE & YES & YES \\
N & 24,720 & 12,240 \\
R² & 0.455 & 0.557 \\
交互项P值 & 0.031** & \\
\end{longtable}
}

替代切点检验进一步表明，以2021为切点或剔除2020过渡年后再比较前后阶段，组间差异均不稳定，因此本文不再将2020或2021作为主文中的首选断点。相关结果见附录。

\subsection{（二）信息中介活跃度}\label{ux4e8cux4fe1ux606fux4e2dux4ecbux6d3bux8dc3ux5ea6}

除制度阶段外，信息中介的活跃程度也可能影响数据利用披露的吸收速度。若分析师跟踪更充分，年报中的相关表述更容易被解释、传播并进入价格，数据利用披露的边际定价价值因而可能更高。

表13第(1)(2)列显示，高分析师关注组中，DU\_kw的系数为-0.0048（t=-4.18），低分析师关注组中为-0.0025（t=-1.75，仅边际显著），交互项P值为0.001。该结果提供了一条补充性边界证据：信息中介越活跃，年报中的数据利用信息越容易被资本市场转化为更快的价格发现。这一证据与ForecastDisp渠道形成互补：前者强调分析师群体对同一企业信息解释更趋一致，后者说明在分析师本身更活跃的企业中，相关披露的定价含量也更容易被释放。

\subsection{（三）上市板块差异}\label{ux4e09ux4e0aux5e02ux677fux5757ux5deeux5f02}

上市板块差异则反映了企业所处信息环境和投资者预期的不同。主板公司经营模式相对成熟、投资者覆盖更广，年报中的数据利用披露未必形成强增量信号；相较之下，非主板企业通常成长性更强、业务不确定性更高，数据利用信息更可能成为投资者重新理解企业价值的重要线索。

表13第(3)(4)列显示，主板公司中DU\_kw的系数为-0.0017（t=-0.98），并不显著；非主板公司中系数为-0.0042（t=-4.01），在1\%水平上显著为负，交互项P值小于0.001。该结果同样提供补充性边界证据：数据要素利用披露的边际定价价值在非主板企业中更强，说明当信息摩擦和成长性不确定性更高时，年报中的相关表述更容易构成有效的信息更新。

\textbf{表13 补充异质性检验}

{\def\LTcaptype{} % do not increment counter
\begin{longtable}[]{@{}lllll@{}}
\toprule\noalign{}
& (1) & (2) & (3) & (4) \\
\midrule\noalign{}
\endhead
\bottomrule\noalign{}
\endlastfoot
维度 & 高分析师关注 & 低分析师关注 & 主板公司 & 非主板公司 \\
因变量 & PriceDelay & PriceDelay & PriceDelay & PriceDelay \\
DU\_kw & -0.0048*** & -0.0025* & -0.0017 & -0.0042*** \\
& (0.0011) & (0.0014) & (0.0017) & (0.0010) \\
控制变量 & YES & YES & YES & YES \\
Firm FE & YES & YES & YES & YES \\
Year FE & YES & YES & YES & YES \\
N & 18,912 & 17,847 & 15,900 & 21,394 \\
交互项P值 & 0.001*** & & \textless0.001*** & \\
\end{longtable}
}

在统一滞后规格下，数字经济核心产业口径的交互项P值为0.507，未能提供稳定的组间差异证据；因此，本文不再将其作为主文异质性的正式结论，仅在附录中作为补充结果简要报告。

\subsection{（四）行业披露水平的同群溢出效应}\label{ux56dbux884cux4e1aux62abux9732ux6c34ux5e73ux7684ux540cux7fa4ux6ea2ux51faux6548ux5e94}

前三条异质性证据说明数据要素利用披露的定价效应受制度阶段、信息中介活跃度和上市板块的调节。一个互补的问题是：除企业自身披露之外，所处行业的整体披露水平是否独立地影响企业股价的信息吸收速度。若行业整体披露水平更高，同业之间的可比性和信息解读规范可能更成熟，投资者在评估单个企业时所面临的解释成本也相应下降，从而形成``同群溢出''。

为此，本文构造行业-年份层面的披露渗透率指标IndPen\_mean，定义为企业所在二级行业在当年的DU\_kw均值，并采用留一法剔除企业自身。将该指标与企业自身DU\_kw同时放入主回归，可分别识别``自身披露''与``行业信息环境''两条渠道的效应。为使系数在量纲上可比，两项指标均作标准化处理。控制变量、固定效应结构和聚类方式与主文一致。

表14第(1)列显示，单独纳入企业自身DU\_kw\_z时，系数为-0.0086（t=-5.10），与主文基准一致。第(2)列同时纳入IndPen\_mean\_z后，企业自身DU\_kw\_z仍显著为负（-0.0067，t=-5.26），而行业披露渗透率IndPen\_mean\_z的系数为-0.0087（t=-2.47），在5\%水平上显著。第(3)列进一步在滞后规格下重复该回归，结果显示行业渗透率的系数增大至-0.0158（t=-4.46），在1\%水平上高度显著，绝对值甚至超过企业自身披露的系数（-0.0070，t=-3.55）。这表明，行业整体披露水平与企业股价延迟之间存在独立且稳健的负向关联，且在时间顺序更清晰的滞后规格下更为突出。

\textbf{表14 行业披露水平的同群溢出效应}

{\def\LTcaptype{} % do not increment counter
\begin{longtable}[]{@{}llll@{}}
\toprule\noalign{}
& (1) & (2) & (3) \\
\midrule\noalign{}
\endhead
\bottomrule\noalign{}
\endlastfoot
规格 & 基准 & 同期 & 滞后 \\
因变量 & PriceDelay & PriceDelay & PriceDelay \\
DU\_kw\_z & -0.0086*** & -0.0067*** & \\
& (0.0017) & (0.0013) & \\
DU\_kw\_z\_lag & & & -0.0070*** \\
& & & (0.0020) \\
IndPen\_mean\_z & & -0.0087** & \\
& & (0.0035) & \\
IndPen\_mean\_z\_lag & & & -0.0158*** \\
& & & (0.0035) \\
控制变量 & YES & YES & YES \\
Firm FE & YES & YES & YES \\
Year FE & YES & YES & YES \\
N & 43,312 & 43,312 & 38,061 \\
\end{longtable}
}

该结果为主文``政策阶段效应''提供了补充性理解。2022年之后，数据要素披露的个体边际效应减弱，但行业整体披露水平仍与更快的价格发现相关。两条证据共同表明：单家企业披露的边际新颖性随制度化推进而下降，但行业层面披露规范的成熟带来了更稳定的信息环境改善。据此，``数据要素利用披露的定价价值''不应被理解为个体层面的单一信号效应，而应理解为``个体披露的相对新颖性''与``行业披露形成的信息共同体''共同作用的结果。需要说明的是，进一步引入企业自身DU\_kw与行业渗透率的交互项时，其系数在两种渗透率度量下方向不一致，本文因此不将其作为主文结论，相关结果见附录。

总体而言，主文异质性证据首先支持``信息新颖性与市场先验预期''这一主要边界条件，并进一步提供``信息中介活跃度''``上市板块信息摩擦''和``行业披露水平的同群溢出''三条补充性边界证据，而非广泛的产业差异叙事。

\input{sections/07_conclusion.tex}

\clearpage
\input{sections/09_references.tex}

\clearpage
\input{sections/appendix.tex}

\end{document}
